%!TEX root = ../../../main.tex
\chapter{Cloud-Abhängigkeit, Lock-in und Funktionsverlust}\label{chap:cloud}
\noindent\textit{Verfasser: Dennis Gerhardt}
\medskip

Beim Smart-TV steht im Vordergrund, welche Daten ein Gerät über seine Nutzer
preisgibt. Bei vielen anderen Haushaltsgeräten reicht die Bindung an den
Hersteller weiter: Sie entscheidet darüber, ob das Gerät überhaupt
funktioniert. Dieses Kapitel untersucht, warum Hersteller solche Geräte an ihre
Server binden und welche Folgen das für Nutzbarkeit, Lebensdauer und Eigentum
hat. Leitend ist die These, dass der Cloud-Zwang überwiegend aus
Geschäftsmodellen folgt und nicht aus technischer Notwendigkeit. Das
eigentliche Problem ist dabei nicht die Cloud selbst, sondern das Fehlen einer
lokalen Rückfallebene.

\section{Cloud-Zentrierung und ihre ökonomischen Treiber}\label{sec:cloud-treiber}

\subsection{Das Gerät als Endpunkt}\label{sec:cloud-endpunkt}

Hoofnagle, Kesari und Perzanowski bezeichnen vernetzte Konsumgüter als
\enquote{tethered}: Sie bleiben nach dem Kauf dauerhaft mit dem Verkäufer
verbunden und werden dadurch von ihm abhängig \cite[S.~785]{hoofnagle:2019}.
Getragen wird diese Bindung von eingebetteter Software, Sensorik und einer
dauerhaften Netzwerkverbindung \cite[S.~789--792]{hoofnagle:2019}. Technisch
bedeutet das, dass die Steuerungslogik nicht mehr im Gerät, sondern in den
Rechenzentren des Herstellers liegt. Das Gerät wird zur Schnittstelle, die
Befehle über Konto, App und Server empfängt, und verarbeitet ohne diese
Verbindung oft nur noch einen Bruchteil seiner Funktionen.

Die Bindung hat durchaus Vorzüge. Hersteller können vernetzte Geräte nach dem
Kauf fortlaufend aktualisieren, verbessern und um neue Funktionen erweitern;
manche Nutzer geben Kontrolle bereitwillig gegen eine ausgereifte Bedienung ab
\cite[S.~804\,f., 808]{hoofnagle:2019}. Diese Vorteile setzen jedoch keine
ausschließlich serverseitige Steuerung voraus. Updates lassen sich ebenso an
Geräte verteilen, deren Grundfunktionen lokal laufen.

Für die Nutzer entsteht aus der Cloud-Zentrierung eine doppelte Abhängigkeit:
von der Verfügbarkeit der Infrastruktur und vom Fortbestand des Anbieters. Beim
Saugroboter betrifft sie gerade die Funktionen, die ihn von einem einfachen
Gerät unterscheiden. Er erfasst die Wohnung mit LiDAR- oder Kamerasensoren und
erzeugt daraus eine Karte; raumgenaue Reinigung, Sperrzonen und Zeitpläne
werden jedoch häufig über App und Herstellerkonto verwaltet. Wie weit diese
Abhängigkeit reicht, zeigt der in Abschnitt~\ref{sec:cloud-neato} untersuchte
Fall Neato. Technisch zwingend ist diese Architektur nicht, denn Karten und
Zonen ließen sich ebenso auf dem Gerät oder im Heimnetz speichern. Sie ist eine
Designentscheidung, deren Gründe vor allem ökonomischer Natur sind.

\subsection{Vendor-Lock-in durch geschlossene Schnittstellen}\label{sec:cloud-lockin}

Die Cloud-Zentrierung schafft die Voraussetzung für einen Vendor-Lock-in. Nach
Hoofnagle et al.\ steigen mit der Bindung der Geräte die Wechselkosten, was
Anbieter strategisch nutzen können, um Wettbewerb zu begrenzen und ihre
Gewinne zu erhöhen \cite[S.~787, 809, 839\,ff.]{hoofnagle:2019}. Eine Studie im
Auftrag des Umweltbundesamts (\acrshort{UBA}) beobachtet entsprechende
Lock-in-Effekte, wenn Geräte nur mit Produkten desselben Herstellers
zuverlässig zusammenarbeiten \cite[S.~16]{jaegererben:2023}. Sie beschreibt
zudem sogenannte Gatekeeper, die rechtliche, finanzielle oder organisatorische
Kontrolle über einzelne Schichten eines Produktsystems ausüben. Solche Akteure
entstehen meist durch Plattform- und Netzwerkeffekte und sind oft Monopolisten
oder Oligopolisten \cite[S.~55]{jaegererben:2023}.

Im Smart Home wird dieser Hebel gezielt eingesetzt. Der Garagentorhersteller
Chamberlain sperrte 2023 Drittanbieter aus seinem myQ-System aus und begründete
dies damit, eine nicht autorisierte Nutzung seines Ökosystems durch fremde Apps
verhindern zu wollen. Die Open-Source-Plattform Home Assistant musste ihre seit
2017 bestehende Integration entfernen, während Partnerunternehmen gegen
Bezahlung weiterhin Zugang erhielten \cite{homeassistant:2023}. Als nicht
autorisiert gilt damit nicht der Zugriff Fremder, sondern die Steuerung des
eigenen Garagentors über eine Software eigener Wahl. Bepreist wird keine neue
Leistung, sondern die Möglichkeit, das eigene Gerät außerhalb des
Hersteller-Ökosystems zu nutzen. Wer sein Garagentor in eine
herstellerübergreifende Steuerung einbinden will, muss auf einen zahlenden
Partner ausweichen oder das Gerät austauschen.

Überträgt man die Bitkom-Definition, nach der digital souveräne Akteure
\enquote{selbstbestimmt und selbstbewusst zwischen Alternativen} wählen können
\cite[S.~7]{bitkom:2015}, auf Verbraucher, so ist vor allem die
Entscheidungssouveränität betroffen: Eine echte Wahl besteht nur noch beim
Kauf, nicht mehr während der Nutzung. Selbst diese Wahl ist eingeschränkt, denn
ob ein Hersteller seine Schnittstellen später schließt, lässt sich beim Kauf
nicht absehen.

\subsection{Hardware as a Service}\label{sec:cloud-haas}

Hinter dem Lock-in steht ein Finanzierungsproblem: Das Gerät wird einmal
verkauft, der Serverbetrieb verursacht dauerhaft Kosten. Abonnements
versprechen den Herstellern daher stetigere Einnahmen
\cite[S.~806]{hoofnagle:2019}. Das Gerät wird so zur Plattform für
wiederkehrende Erlöse, zu denen auch die im vorigen Kapitel beschriebene
Verwertung von Nutzungsdaten zählt. Hoofnagle et al.\ beschreiben, wie sich
dadurch der Charakter des Geschäfts verschiebt: Was wie ein Kauf aussieht, kann
sich in ein Dienstverhältnis verwandeln, dessen Bedingungen der Anbieter
nachträglich ändert. Statt ein Gerät auf einen Schlag abzuschalten, können
Hersteller zudem Funktionen schrittweise entfernen oder die Leistung
verschlechtern \cite[S.~814\,f.]{hoofnagle:2019}.

Wie weit dieses Modell reichen kann, zeigt der norwegische Hersteller
Futurehome. Er hatte seine Smart-Home-Zentrale an rund 38\,000 Kunden mit dem
Versprechen verkauft, dass nach dem Kauf keine laufenden Gebühren anfallen.
Nach der Insolvenz 2025 verlangte die Nachfolgegesellschaft eine Jahresgebühr;
ohne sie entfallen App-Steuerung, Automationen und Energiefunktionen, und die
Geräte lassen sich nur noch manuell bedienen. Durchgesetzt wurde dies per
Firmware-Update auf bereits bezahlten Zentralen \cite{schoutsen:2025}. Der Fall
belegt die Leitthese besonders deutlich: Die Funktionen liefen zuvor auf der
Zentrale im Haushalt und wurden nicht aus technischen, sondern aus
wirtschaftlichen Gründen gesperrt. Die Kunden standen dabei vor einer
Zwangslage: Wer nicht zahlte, verlor den Nutzen der Zentrale und der
eingebundenen Geräte; wer zahlte, war auf ein Unternehmen angewiesen, dessen
Vorgänger sein Versprechen nicht hatte halten können.

Abonnements sind nicht per se illegitim. Serverkosten sind real, und
Hoofnagle et al.\ sehen in nutzungsabhängigen Modellen sogar die Chance, Käufern
den Zugang zu erleichtern, die sich das vollständige Produkt nicht leisten
könnten \cite[S.~807]{hoofnagle:2019}. Ein beim Kauf ausgewiesenes Abonnement
erlaubt zudem eine informierte Entscheidung. Problematisch ist die nachträgliche
Bepreisung bereits bezahlter Funktionen. Dass der Markt dies nicht korrigiert,
erklärt sich aus einer Informationsasymmetrie: Vernetzte Produkte werden durch
die dauerhafte Bindung von Such- zu Erfahrungsgütern, deren Eigenschaften sich
erst nach dem Kauf zeigen \cite[S.~809]{hoofnagle:2019}. Wie lange ein Dienst
betrieben wird und welche Funktionen später kostenpflichtig werden, können
Käufer nicht erkennen. Hoofnagle et al.\ schlagen deshalb vor, Hersteller zur
Angabe der erwarteten Lebensdauer und der Obsoleszenzrisiken zu verpflichten
\cite[S.~866]{hoofnagle:2019}.

\section{Gefährdung von Lebensdauer und Funktionsfähigkeit}\label{sec:cloud-folgen}

Die materiellen Folgen dieser Mechanismen zeigen sich am deutlichsten, wenn
ein Hersteller seine Dienste einstellt oder nachträglich einschränkt.

\subsection{Fallbeispiel Saugroboter: Die Abschaltung der Neato-Cloud}\label{sec:cloud-neato}

Vorwerk hatte die 2017 übernommene Tochter Neato Robotics 2023 geschlossen und
zugesagt, die Cloud-Dienste der Saugroboter mindestens fünf Jahre
weiterzubetreiben. Nur zwei Jahre später kündigte das Unternehmen die
Abschaltung an und verwies auf gestiegene Anforderungen an die Cybersicherheit,
denen die bestehende Infrastruktur nicht mehr genüge. Seitdem lassen sich die
Roboter nur noch per Taste starten und reinigen dann die gesamte Wohnung in
einem Durchgang; zuvor angelernte Räume und Sperrzonen erkennen sie
voraussichtlich nicht mehr. Neben der Cybersicherheit nannte Vorwerk auch
wirtschaftliche Gründe: Neato habe seine Ziele verfehlt, und der Konzern ziehe
sich aus nicht wachsenden Geschäftsfeldern zurück \cite{heise:2025neato}.

Gerade die Funktionen, die einen Saugroboter von einem einfachen Haushaltsgerät
unterscheiden, hingen also an der Cloud. Dass die angelernten Räume mit ihr
verschwinden, zeigt zudem, dass die in der Wohnung erhobenen Kartendaten
faktisch nicht beim Nutzer, sondern beim Hersteller lagen. Aus
Unternehmenssicht ist die Abschaltung folgerichtig, denn eine Cloud für ein
eingestelltes Produkt verursacht Kosten ohne Erlöse. Diese Kosten werden jedoch
auf die Käufer verlagert, deren Geräte an Wert verlieren.

Der Fall zeigt zweierlei. Erstens war die Zusage nicht belastbar: Die
Abschaltung kam rund drei Jahre früher als versprochen. Zweitens bewegte erst
behördlicher Druck den Hersteller. Die italienische Wettbewerbsbehörde
\acrshort{AGCM} leitete nach der Abschaltung im April~2026 ein Verfahren ein und
erklärte im November Zusagen der Vorwerk-Gruppe für verbindlich: Wer seit
Oktober~2020 einen Neato-Roboter gekauft hat, erhält europaweit gegen Rückgabe
ein Vorwerk-Neugerät oder einen Gutschein über 500~Euro; für das Modell VR7
gilt zudem eine zehnjährige Garantie auf digitale Dienste einschließlich
Sicherheitsupdates \cite{agcm:2026b}. Die Entschädigung besteht allerdings im
Umstieg auf ein Gerät desselben Konzerns, das ursprüngliche Gerät bleibt
eingeschränkt.

\subsection{Der Cloud-Tod: Gigaset Smart Home}\label{sec:cloud-gigaset}

Im Extremfall wird funktionsfähige Hardware vollständig wertlos. Im
Insolvenzverfahren der Gigaset Communications GmbH übernahm der Käufer der
übrigen Unternehmensteile die Sparte Smart Home/Care nicht. Die Cloud-Dienste
wurden deshalb zum 29.~März~2024 eingestellt; Apps, Sensoren und vernetzte
Geräte waren danach nicht mehr nutzbar, gebuchte Kamera-Servicepakete endeten,
und nur der lokale Alarm des Rauchmelders funktionierte weiter. Eine Rückgabe
der Geräte war ausgeschlossen; Ansprüche konnten Kunden lediglich im
Insolvenzverfahren anmelden \cite{gigaset:2024a}.

Betroffen waren Geräte, die gerade zum Schutz der Wohnung angeschafft worden
waren. Bei Sensoren und Kameras ist ein solcher Ausfall nicht nur ein
Komfortverlust, sondern eine Schutzlücke, die Nutzer womöglich erst bemerken,
wenn eine Meldung ausbleibt. Über den Fortbestand eines Dienstes entscheidet in der Insolvenz die
Verwertung der Masse, nicht der Wille des Herstellers. Hoofnagle et al.\ führen
das Scheitern von Unternehmen deshalb als eigenständiges Risiko vernetzter
Produkte auf, das Käufern wertlose Geräte hinterlassen kann
\cite[S.~816]{hoofnagle:2019}. Im engen Sinne Tusikovs, die unter Bricking eine
vom Hersteller veranlasste Störung der Software mit dem Ziel versteht, die
Funktion eines Produkts zu beeinträchtigen \cite[S.~2]{tusikov:2019}, handelt
es sich hier zwar nicht um Bricking, da Gigaset die Geräte nicht unbrauchbar
machen wollte. Für die Käufer ist das Ergebnis jedoch dasselbe.

Dass auch wirtschaftlich starke Unternehmen Geräte bewusst stilllegen, zeigte
bereits 2016 der Fall Revolv: Nach der Übernahme des Herstellers stellte Nest
den Verkauf der Smart-Home-Zentrale ein und machte später sämtliche Geräte
funktionsunfähig \cite[S.~811\,f.]{hoofnagle:2019}. Weniger drastisch verlief
die Abschaltung bei Belkin. Der Hersteller kündigte 2025 an, zum
31.~Januar~2026 Cloud-Dienste und App für die meisten Wemo-Geräte einzustellen,
darunter auch Geräte, die erst vor weniger als zwei Jahren gekauft worden
waren. Bereits eingerichtete Geräte lassen sich über eine lokale Schnittstelle
etwa in Home Assistant weiter steuern, neu einrichten lassen sie sich ohne
Herstellercloud jedoch nicht \cite{schoutsen:2025}. Eine lokale Schnittstelle
mildert den Funktionsverlust also, beseitigt ihn aber nicht.

\subsection{Softwarebedingte Obsoleszenz und prekäres Eigentum}\label{sec:cloud-obsoleszenz}

Die UBA-Studie fasst unter Softwareobsoleszenz Phänomene, bei denen die
Qualität einer Software so nachlässt, dass ein Produktsystem seine Funktion
für eine erhebliche Zahl von Nutzern nicht mehr angemessen erfüllt
\cite[S.~13]{jaegererben:2023}. Die Ursachen sieht sie weniger in der Technik
als in wirtschaftlichen und organisatorischen Bedingungen wie eingestelltem
Support \cite[S.~5]{jaegererben:2023}; Cloud-Dienste behandelt sie als eigene
Kategorie externer Abhängigkeit \cite[S.~55]{jaegererben:2023}. Bei der
Auswertung von Nutzerforen und Bewertungen stieß sie auf smarte Geräte, die nur
in Verbindung mit Server-Infrastruktur funktionieren und teils schon kurz nach
dem Kauf unbrauchbar wurden \cite[S.~16]{jaegererben:2023}. Unter den Begriff
fallen Abschaltungen wie bei Neato und Gigaset ebenso wie einseitige
Verschlechterungen per Update wie bei Futurehome.

Dahinter steht ein verändertes Verständnis von Eigentum. Tusikov beschreibt
den Kauf vernetzter Geräte als prekäre Form des Eigentums: Käufer besitzen die
Hardware, nutzen die Software aber nur im Rahmen einer Lizenz, deren
Bedingungen der Hersteller nach dem Kauf ändern kann. Über die Software
kontrolliert er so auch die Hardware \cite[S.~2, 11]{tusikov:2019}. Hoofnagle
et al.\ schlagen daher vor, dass Verbraucher die Bindung an den Hersteller
lösen und ihr Gerät einem anderen Anbieter zuführen können, sobald die
Unterstützung endet oder der Hersteller ausfällt \cite[S.~868]{hoofnagle:2019}.
Tabelle~\ref{tab:cloud-faelle} fasst die betrachteten Fälle zusammen.

\begin{table}[htbp]
  \centering
  \caption{Abschaltungen und nachträgliche Eingriffe bei vernetzten Haushaltsgeräten}
  \label{tab:cloud-faelle}
  \footnotesize
  \setlength{\tabcolsep}{4pt}
  \renewcommand{\arraystretch}{1.25}
  \begin{tabularx}{\textwidth}{@{}>{\raggedright\arraybackslash}p{2.7cm}
      >{\raggedright\arraybackslash}p{3.3cm}
      >{\raggedright\arraybackslash}X
      >{\raggedright\arraybackslash}p{2.9cm}@{}}
    \toprule
    \textbf{Fall (Jahr)} & \textbf{Anlass} & \textbf{Verbleibende Funktion}
      & \textbf{Betroffene Dimension} \\
    \midrule
    Revolv (2016) & Übernahme durch Nest
      & Keine; Zentrale funktionsunfähig & Anwendung \\
    Chamberlain myQ (2023) & Sperre für Drittanbieter
      & Steuerung nur über Hersteller-App und zahlende Partner
      & Entscheidung \\
    Gigaset Smart Home (2024) & Insolvenz, Sparte nicht übernommen
      & Nur lokaler Alarm des Rauchmelders & Anwendung \\
    Futurehome (2025) & Insolvenz, Abo-Zwang
      & Manuelle Bedienung der Einzelgeräte ohne App und Automationen
      & Entscheidung, Anwendung \\
    Neato (2025) & Konzernentscheidung
      & Start per Taste, Reinigung ohne Räume und Sperrzonen
      & Anwendung, Daten \\
    Belkin Wemo (2026) & Einstellung von Cloud und App
      & Bereits eingerichtete Geräte über lokale Schnittstelle
      & Anwendung \\
    \bottomrule
  \end{tabularx}
  \par\smallskip
  \raggedright\scriptsize
  Quelle: Eigene Darstellung auf Basis von
  \cite[S.~811\,f.]{hoofnagle:2019} sowie
  \cite{homeassistant:2023,gigaset:2024a,schoutsen:2025,heise:2025neato}.
\end{table}

\section{Zwischenfazit}\label{sec:cloud-zwischenfazit}

Die Cloud-Abhängigkeit vernetzter Haushaltsgeräte berührt alle drei
Souveränitätsdimensionen: die Anwendungssouveränität, wenn die Nutzbarkeit am
Fortbestand eines Anbieters hängt, die Entscheidungssouveränität, wenn
Hersteller Funktionen nach dem Kauf sperren oder bepreisen, und die
Datensouveränität, wenn Wohnungskarten nur beim Hersteller vorliegen. Die Fälle
bestätigen die Leitthese, dass lokale Funktionen technisch möglich sind und aus
wirtschaftlichen Gründen fehlen oder entfernt werden. Daraus folgt keine
Forderung nach Autarkie, die auch der Bitkom ablehnt \cite[S.~7]{bitkom:2015},
wohl aber nach einer lokalen Rückfallebene für Kernfunktionen, offenen
Schnittstellen und einem geregelten Ausstiegspfad beim Supportende, wie sie in
Teilen auch die UBA-Studie empfiehlt \cite[S.~18\,f.]{jaegererben:2023}.
Inwieweit lokale Zentralen, offene Standards und alternative Firmware dies
leisten, untersuchen die folgenden Kapitel.
